\subsection{Resolution sensitivity design}\label{sec:resolution-design}

The resolution study varies the budgets of the primary particle estimator while holding the data-generating process fixed at the four income regimes IC0 to IC3, at $n=1000$, and on seeds $0,\ldots,9$. Grid resolution is varied over $K\in\{5,25,49\}$ at $M=10$ and particle resolution over $M\in\{5,10,25\}$ at $K=25$. The one-factor slices are run first, each against the incumbent pair $(K,M)=(25,10)$. A full $3\times3$ factorial extension follows only if the one-factor slices show a material change, and it is allocated only to IC1 and IC3, the regime with a policy effect and the regime with the most difficult overlap. The forest comparators do not use $M$, so their resolution rows are repeated once per $K$ and are never duplicated across $M$.

The retention rule was fixed before the outcomes were inspected and is evaluated on the native-grid metrics of each cell. Write $R^{t}_{r}(K,M)$ for the between-seed mean WCF error on target $t\in\{\mathrm{ref},\mathrm{law}\}$ (reference-TCATE and conditional law) in regime $r$, and $\mathrm{SE}^{t}_{r}(K',M')$ for the between-seed standard error of the paired difference between $(K',M')$ and $(25,10)$. The incumbent pair $(25,10)$ is retained when every one-factor alternative $(K',M')\in\{(5,10),(49,10),(25,5),(25,25)\}$ satisfies
\begin{equation}\label{eq:sens-retention-rule}
\frac{\lvert R^{t}_{r}(K',M')-R^{t}_{r}(25,10)\rvert}{R^{t}_{r}(25,10)}\le 0.10,
\qquad t\in\{\mathrm{ref},\mathrm{law}\},\quad r\in\{\mathrm{IC1},\mathrm{IC3}\},
\end{equation}
and when no larger grid or particle count with $K'>25$ or $M'>25$ improves both targets by more than $2\,\mathrm{SE}^{t}_{r}(K',M')$ in both regimes. If the rule fails, the study reports the selected pair and the complete sensitivity table rather than silently changing the primary specification.

Because native-grid error targets change their estimand with $K$, the study also evaluates every fitted law on a common dense target. The common level set is
\begin{equation}\label{eq:sens-common-grid}
u_j=\frac{j-1/2}{199},\qquad j=1,\ldots,199,
\end{equation}
and particle coordinates and forest atoms are interpolated linearly in the quantile level, with constant continuation from the nearest endpoint outside the fitted grid, so no predicted value is extrapolated beyond the support of its own grid. The true data-generating curves and the regime reference vector are constructed directly on these levels from the regime specification, never interpolated from their coarse versions, and every method within a replication is evaluated against the same dense truth, test covariates, and evaluation bandwidth. A second interior level set $0.1+0.8u_j$ restricts the diagnostic to levels on which none of the proposed values $K\in\{5,25,49\}$ requires endpoint continuation, which separates tail coverage from the interior fit. Native-grid reference targets are regenerated on each cell's own $K$-grid; native, common, and interior targets are recorded under distinct identifiers and are never pooled.

The dense-grid calibrated functional effects use the interpolated particle clouds for the out-of-fold nuisance means and the interpolation of the observed coarse quantile curve for the observed functional values. This is a distinct information setting from a dense observed curve: the calibrated contrast targets the dense estimand while the observations remain coarse, so the augmentation identity is not exact at the dense level, and the interior reference targets inherit the same distinction. Native-grid calibrated effects are never relabelled as dense-grid effects.

The retained or selected pair is confirmed on fresh seeds $20,\ldots,39$ for $K\in\{25,49\}$ at $M=10$ over IC0 to IC3, with no reselection on those seeds. The confirmation is a stability check on twenty replications, not an equivalence test. Wall time and process peak memory are recorded for every cell. The study used seven concurrent single-thread workers, each pinned to one numerical thread, and the factorial extension was allocated only to IC1 and IC3.

\subsection{Assignment mechanisms and propensity models}\label{sec:assignment-design}

A second design family, labelled SYM, uses six independent $\mathrm{Unif}[-1,1]$ covariates and the generative form of equation~\eqref{eq:dgp}, but with a strictly symmetric inner law. Its outcome surfaces are
\begin{align*}
\mu(x)&=0.70\sin(\pi x_1)+0.35x_2x_3+0.20\cos(\pi x_4),\\
\tau(x)&=0.25+0.15\sin(\pi x_2),\\
s(x)&=0.12+0.08x_5^2-0.05x_3,
\end{align*}
with location $m_a(x)=\mu(x)+a\tau(x)$, log-scale $s_a(x)=s(x)$, shape coefficient $\gamma_a(x)=0$, no outer log-scale shock ($\eta_a\equiv 0$), and an outer location shock $\xi_a\sim N(0,0.25^2)$. Since $\psi(z;0)=z$, the arm law reduces to $q_k(Y^a)=\mu(X)+a\tau(X)+\xi_a+\exp\{s(X)\}z_k$, so the conditional quantile contour is pointwise symmetric in both arms: each coordinate $Q_k^a$ is symmetrically distributed about $\mu(X)+a\tau(X)+\exp\{s(X)\}z_k$, and $x_6$ enters no outcome surface.

Four assignment mechanisms share this outcome law. SYM-RANDOM uses the constant propensity $e(x)=0.5$. SYM-LIN uses the logistic link of a linear index, $e(x)=C_{.05,.95}\{\sigma(1.4x_1-1.1x_2+0.4x_3)\}$. SYM-NL uses a nonlinear propensity surface, $e(x)=C_{.05,.95}\{\sigma(1.6\sin(\pi x_1)+0.8x_2x_3-0.7x_4^2+0.35x_5)\}$. SYM-MU makes assignment a function of the baseline outcome location, $e(x)=C_{.05,.95}\{\sigma(1.2\mu(x)+0.5x_5)\}$. In all four cases the propensity is a function of pre-treatment covariates only, not of the potential outcomes; SYM-MU uses the baseline surface $\mu$ rather than the arm-specific law. Each mechanism has a null companion obtained by setting $\tau(x)=0$. The four mechanisms are fitted at $K=25$, $M=10$, ten seeds, and sample sizes $n\in\{500,1000\}$. The principal comparisons under SYM-LIN, SYM-NL, and SYM-MU are against the distributional forest comparators \citep{DRF-paper,naf2026causaldrf}, with SYM-RANDOM as a calibration control.

The alignment pair keeps the same outcome law and contrasts two nonlinear propensity indices with the same marginal law. SYM-ALIGN uses $e(x)=\sigma\{1.6\sin(\pi x_1)\}$ and SYM-IRREL uses $e(x)=\sigma\{1.6\sin(\pi x_6)\}$, where $x_6$ enters no outcome surface. Because the covariates are independent and identically distributed, the two propensity variables have the same marginal distribution by construction, so the pair isolates the relation between assignment and prognosis rather than the propensity spread. The alignment pair uses the same budgets and sample sizes as the four mechanisms, and null companions are included for both members.

Design checks on $200{,}000$ draws per regime confirm the intended overlap and balance before any estimator score is interpreted. Under SYM-NL the propensity ranges from $0.050$ to $0.932$ with treated fraction $0.4557$, and the largest standardized mean differences are $1.131$ for $\mu(X)$, $1.092$ for $\sin(\pi x_1)$, $0.819$ for $x_1$, and $0.212$ for $x_2x_3$. Under SYM-MU the propensity ranges from $0.131$ to $0.870$ with treated fraction $0.4992$, and the largest standardized mean differences are $0.615$ for $\mu(X)$, $0.574$ for $\sin(\pi x_1)$, $0.444$ for $x_1$, and $0.258$ for $x_5$. The nonlinear prognostic features carry the largest imbalances, which raw covariate means alone would not reveal. The generated inner laws pass the pointwise symmetry and monotonicity checks, with grid skewness at numerical precision (maximum absolute value $1.8\times10^{-15}$ across arms and regimes). The SYM-ALIGN and SYM-IRREL propensity marginals agree to a maximum absolute difference of $0.0018$.

The AIPW calibration layer accepts a binary-probability estimator factory that is fitted inside each training fold and returns probabilities on the same clipped scale. The default factory is the five-fold logistic regression with clipping to $[0.02,0.98]$ used in the primary fits. The propensity sensitivity adds a random-forest factory (200 trees, minimum leaf size 5, single-threaded), a gradient-boosting factory at the library defaults, and an oracle diagnostic that substitutes the true propensity; the oracle is a diagnostic benchmark and is explicitly not a feasible competitor. All sensitivity fits use three outcome folds and contrast candidates $\{0,50,500\}$, holding the outcome learner, $K$, $M$, seeds, and test designs fixed. The random-forest and gradient-boosting factories use no held-out tuning beyond their fixed settings, so their contrasts measure the propensity model family rather than a tuned alternative, and the comparison under nonlinear assignment is a misspecification stress test that does not by itself establish the doubly robust property.
